\subsection{第五係数零表示の二族と二固定表示}
\label{subsec:cubic-as4-zero-fifth}
引き続き\cite{AS4lowpinned}の具体表示だけを参照する。
第五係数零の開部分FZ2--FZ3と境界FB2にある次の四表示を扱う。
この節で $M$ と書く族は、前節の $M_c$ とは異なる。

\begin{proposition}\label{prop:cubic-zero-fifth-m}
$\lambda=1$, $p\ne1$, $\varepsilon=0,1$ とし、
\[
 x_0=x_1=1,\quad x_2=p,\quad x_9=\varepsilon,\quad
 x_{14}=\frac1{16(p-1)},\quad x_{19}=-\frac1{16(p-1)},\quad
 x_{20}=\frac{p}{16(p-1)}
\]
として他の $x_j$ を零にした表示 $M^\varepsilon_p$ は、全値でH3包絡にならない。
\end{proposition}
\begin{proof}
全関係moduleから得る三次・四次商座標の分母は $p-1$ の冪だけである。
これらは全関係を消し自由列上でidentityなので、AS4実現の各 $p\ne1$ で
次数別次元20,35の真の商基底となる。
中心三次行列の20次minorは、それぞれ
\[
 \frac{3p+1}{2^{33}(p-1)^{11}},\qquad
 \frac{(p+3)(3p+1)}{2^{23}(p-1)^{13}}
\]
である。残る $p=-1/3$ と、$\varepsilon=1$ における $p=-3$ では、
$\mathbb Q$ 上で関係行列を別に構成して非零20次minorを確認した。
従って全値で中心三次は零である。
また全 $p$ の最高tensor $R A R^{\mathsf T}$ の非零反巡回性を独立検算した。
AS4実現では $\mu=I$ だから、中心三次の消滅がH3を排除する。
特に $p=0$ を除外していない。
\end{proof}

\begin{proposition}\label{prop:cubic-zero-fifth-fixed}
次の二固定表示 $N,C$ もH3包絡にならない。
\[
\begin{array}{c|l}
 N&\lambda=2,\quad \zeta^2=-3,\quad x_0=x_1=1,\quad x_2=-2,\\
 &x_8=\zeta/4,\quad x_{10}=x_{11}=x_{16}=\zeta/6,\quad x_{15}=-\zeta/12,\\
 &x_{14}=x_{19}=x_{20}=1/12;\\\hline
 C&\lambda=1,\quad x_0=x_1=x_2=1,\\
 &x_{10}=x_{16}=i/8,\quad x_{11}=x_{15}=-i/8,\\
 &x_{14}=x_{19}=x_{20}=1/64.
\end{array}
\]
記載のない $x_j$ は零とする。
\end{proposition}
\begin{proof}
$N$ は $\mathbb Q(\sqrt{-3})$、$C$ は $\mathbb Q(i)$ 上で
三次・四次商座標と中心三次行列を構成し、非零20次minorを確認した。
係数共役はそれぞれ $\zeta$ または $i$ の他方の符号も覆う。
両方の非零反巡回最高tensorも全係数で確認したので、AS4実現なら $\mu=I$。
従って中心三次の消滅が主張を与える。
\end{proof}

全値の分母支持・例外minor・数体計算は
\texttt{verify\_\allowbreak{}cubic\_\allowbreak{}as4\_\allowbreak{}four\_\allowbreak{}simple\_\allowbreak{}zero\_\allowbreak{}fifth.py}
で再現できる。AS4存在十分性・全同型・全射影型被覆は輸入していない。
第五係数零境界の別表示 $R^+_1,R^-_1$ は、この節の四表示には含めない。
